%!TEX root = ../thesis.tex
% створюємо перелік умовних позначень, скорочень і термінів
\shortings
    % ФТІ --- Фізико-технічний інститут

    % \kafedraShort{} --- \kafedra{}
    % Autoregressive Transformer --- авторегресивний трансформер

    % baseline model --- базова модель

    % cGAN --- умовна генеративно-змагальна мережа (англ. conditional Generative Adversarial Network)

    % CLIP --- контрастивне попереднє навчання мови та зображень (англ. Contrastive Language-Image Pre-training)

    % CNN --- згорткова нейронна мережа (англ. Convolutional Neural Network)

    % Diffusion Model --- дифузійна модель

    % FID --- відстань Фреше (англ. Fréchet Inception Distance)

    % GAN --- генеративно-змагальна мережа (англ. Generative Adversarial Network)

    % KID --- ядерна відстань (англ. Kernel Inception Distance)

    % LPIPS --- перцептивна метрика на базі глибоких ознак (англ. Learned Perceptual Image Patch Similarity)

    % MAE --- середньоабсолютна похибка (англ. Mean Absolute Error)

    % MSE --- середньоквадратична похибка (англ. Mean Squared Error)

    % PSNR --- пікове відношення сигналу до шуму (англ. Peak Signal-to-Noise Ratio)

    % SSIM --- індекс структурної схожості (англ. Structural Similarity Index)

    % SSM --- селективна модель простору станів (англ. Selective State Space Model)

    % Transformer --- трансформер


    % VRAM --- відеопам'ять (англ. Video Random Access Memory)
    \NewTblrTheme{abbrvlist}{
        \SetTblrTemplate{caption}{empty}
        \SetTblrTemplate{capcont}{empty}
        \SetTblrTemplate{contfoot}{empty}
        \SetTblrTemplate{conthead}{empty}
    }

    \noindent\textbf{Скорочення}
    \vspace{0.5em}

    \begin{longtblr}[
        theme = abbrvlist 
    ]{
        width = \textwidth,
        colspec = {l @{\quad---~~} X[l]}, 
        % rowsep = 0.3em, 
    }
        cGAN & умовна генеративно-змагальна мережа (англ. conditional Generative Adversarial Network) \\
        CLIP & контрастивне попереднє навчання мови та зображень (англ. Contrastive Language-Image Pre-training) \\
        CNN & згорткова нейронна мережа (англ. Convolutional Neural Network) \\
        FID & відстань Фреше (англ. Fréchet Inception Distance) \\
        GAN & генеративно-змагальна мережа (англ. Generative Adversarial Network) \\
        KID & ядерна відстань (англ. Kernel Inception Distance) \\
        LPIPS & перцептивна метрика на базі глибоких ознак (англ. Learned Perceptual Image Patch Similarity) \\
        MAE & середньоабсолютна похибка (англ. Mean Absolute Error) \\
        MSE & середньоквадратична похибка (англ. Mean Squared Error) \\
        PSNR & пікове відношення сигналу до шуму (англ. Peak Signal-to-Noise Ratio) \\
        SSIM & індекс структурної схожості (англ. Structural Similarity Index) \\
        SSM & селективна модель простору станів (англ. Selective State Space Model) \\
        VRAM & відеопам'ять (англ. Video Random Access Memory) \\
    \end{longtblr}

    \vspace{1.5em}
    \noindent\textbf{Терміни}
    \vspace{0.5em}

    \begin{longtblr}[
        theme = abbrvlist 
    ]{
        width = \textwidth,
        colspec = {l @{\quad---~~} X[l]}, 
        rowsep = 0.3em, 
    }
        Autoregressive Transformer & авторегресивний трансформер \\
        baseline model             & базова модель \\
        batch size                 & розмір пакету \\
        Diffusion Model            & дифузійна модель \\
        Transformer                & трансформер \\
        weight decay               & коефіцієнт згасання ваг \\
    \end{longtblr}

    % (БУДЬ ЛАСКА, ПРОСЛІДКУЙТЕ, ЩОБ НОМЕР СТОРІНКИ СПІВПАДАВ ІЗ СПРАВЖНІМ! Це залежить від того, наскільки великим є ваш зміст.
    % Номер сторінки проставляється у файлі thesis.tex, рядок 31.)